\documentclass[11pt,twocolumn]{article}

\usepackage[margin=0.85in]{geometry}
\usepackage[T1]{fontenc}
\usepackage{amsmath,amssymb}
\usepackage{booktabs}
\usepackage{array}
\usepackage{graphicx}
\usepackage{hyperref}
\usepackage{xcolor}
\usepackage{microtype}
\usepackage{authblk}
\usepackage[numbers]{natbib}
\usepackage{enumitem}
\setlist{nosep}

\hypersetup{colorlinks=true,linkcolor=blue!60!black,citecolor=blue!60!black,urlcolor=blue!60!black}

\title{The Standard Genetic Code Is Not Locally Optimal\\Under Simple Chemical Distance}

\author[1]{Jonathan Washburn}
\affil[1]{Independent Researcher}

\date{March 2026}

\begin{document}
\maketitle

\begin{abstract}
We evaluate the local optimality of the standard genetic code under a
simple chemical-distance metric that measures the physicochemical cost of
amino acid substitutions caused by single-nucleotide mutations. Using an
exhaustive pairwise-swap search over all $\binom{20}{2} = 190$ amino-acid
transpositions within the degeneracy-preserving assignment space, we find
that the standard code is \textbf{not} a local minimum: at least 42 pairwise
swaps strictly reduce total weighted mutation cost. The largest improving
swap (Arg\,$\leftrightarrow$\,Gln) reduces cost by 10.26\%. At the same
time, a degeneracy-preserving random baseline of 10,000 shuffled codes shows
that the standard code still outperforms 99.47\% of sampled alternatives
under this metric. The code is therefore globally strong relative to random
codes but not locally optimal in its immediate swap neighborhood. These
results are formalized in Lean~4
with a machine-checked proof that the standard code is not a local minimum.
We conclude
that any serious optimality claim for the genetic code must account for
structure beyond simple physicochemical substitution cost, such as
position-dependent mutation effects, biosynthetic pathway constraints, and
translation-machinery compatibility.
\end{abstract}

\section{Introduction}

The standard genetic code maps 61 sense codons to 20 amino acids via a
famously non-random pattern~\citep{Woese1965, Crick1968}. Since at least
Haig and Hurst~\citep{HaigHurst1991}, a substantial literature has asked
whether this mapping is optimal with respect to minimizing the deleterious
effects of point mutations. The prevailing view, strengthened by the
influential work of Freeland and Hurst~\citep{FreelandHurst1998}, is that
the standard code is near-optimal: it is better than the vast majority of
randomly generated alternatives at minimizing a weighted chemical-distance
objective.

This near-optimality claim, however, has always been stated relative to
random codes. The question of \emph{local} optimality---whether any small
perturbation of the code can strictly improve the objective---has received
less attention. In this paper, we address this question directly.

We define a simple chemical-distance metric based on four amino acid
physicochemical properties (molecular weight, hydrophobicity, charge, and
side-chain volume) and evaluate all 190 pairwise amino-acid-label swaps
within the degeneracy-preserving assignment space. We find that \textbf{42
out of 190 swaps strictly reduce the total weighted mutation cost}. The
standard code is therefore not a local minimum under this metric.

This is not a refutation of the genetic code's optimality. Rather, it is a
precise characterization of what a simple metric can and cannot explain,
and a motivation for identifying the richer objective that the code
actually minimizes.

The central point is slightly subtler, and stronger, than a bare negative
result. Under this chemistry-only objective the standard code performs
extremely well against random degeneracy-preserving shuffles, yet still
admits many improving local perturbations. The natural conclusion is not
that the code lacks structure, but that the wrong objective is being
optimized.

All results are formalized in Lean~4 and verified by the Lean type checker.
The central formal statement is that there exist amino acids $a \neq b$ such
that swapping $a$ and $b$ strictly reduces the total code cost under
chemistry-only weights. The exact theorem and module are given in the
Formal Verification section below. The full formalization is available in the
\texttt{IndisputableMonolith.Genetics} namespace (49 modules, 8,000+ lines,
zero axioms, zero unproved assertions).

\section{Definitions}

\subsection{Chemical Distance}

For amino acids $a_1, a_2$ with physicochemical properties
$(w_i, h_i, c_i, v_i)$ denoting molecular weight (Da), Kyte--Doolittle
hydrophobicity ($\times 10$, integer), charge at pH~7 ($\{-1, 0, +1\}$),
and side-chain volume (\AA$^3$), the chemical distance is
\begin{equation}\label{eq:chemdist}
  d(a_1, a_2) = \Delta w^2 + \Delta h^2 + 100\,\Delta c^2 + \Delta v^2
\end{equation}
where $\Delta w = w_1 - w_2$, etc. The charge term carries a factor of 100
to penalize charge-changing substitutions, which are among the most
biologically disruptive~\citep{Grantham1974}. This metric is symmetric,
non-negative, and zero if and only if $a_1 = a_2$.

\subsection{Mutation Model}

Single-nucleotide mutations are weighted by a transition--transversion
asymmetry: transitions (purine$\leftrightarrow$purine or
pyrimidine$\leftrightarrow$pyrimidine) receive weight~2, transversions
receive weight~1. This reflects the approximately 2:1 transition bias
observed in most genomes~\citep{Wakeley1996}. Same-nucleotide ``mutations''
carry weight~0.

For a codon $c$ with single-nucleotide mutants $\{m_1, \ldots, m_9\}$,
the \textbf{vulnerability} of $c$ under assignment $A$ is
\begin{equation}\label{eq:vuln}
  \begin{aligned}
    V_A(c) &= \sum_{i=1}^{9} \mu(c, m_i)\, K_A(c,m_i), \\
    K_A(c,m_i) &=
    \begin{cases}
      d(A(c), A(m_i)) & \text{if both outputs are sense codons,} \\
      P_{\text{stop}} & \text{if } A(c) \text{ is sense and } A(m_i) \text{ is stop,} \\
      0 & \text{if } A(c) \text{ is stop.}
    \end{cases}
  \end{aligned}
\end{equation}
where $\mu(c, m_i)$ is the transition--transversion weight and
$P_{\text{stop}} = 50{,}000$ is a nonsense penalty.

\subsection{Total Code Cost}

The total cost of an assignment $A$ is the sum of vulnerabilities over all
64~codons:
\begin{equation}\label{eq:total}
  C(A) = \sum_{c \in \text{codons}} V_A(c).
\end{equation}

\subsection{Assignment Space}

An \textbf{assignment} is a bijection $A\colon \text{AminoAcid} \to
\text{AminoAcid}$ that relabels which amino acid each codon family encodes,
while preserving the degeneracy pattern (stop codons remain stop codons,
and all codons in a synonymous family receive the same label). The identity
assignment recovers the standard genetic code.

The \textbf{pairwise swap neighborhood} consists of all
$\binom{20}{2} = 190$ transpositions $(a, b)$ with $a \neq b$, where the
codon families for $a$ and $b$ exchange their amino-acid labels while all
other families remain unchanged.

\subsection{Local Optimality}

The standard code is \textbf{locally optimal} under cost function $C$ if
$C(\text{id}) \leq C(\sigma_{a,b})$ for every transposition $\sigma_{a,b}$
in the swap neighborhood. It is \textbf{not locally optimal} if there exist
$a \neq b$ with $C(\sigma_{a,b}) < C(\text{id})$.

\section{Results}

\subsection{The Standard Code Is Not Locally Optimal}

Under the chemical-distance metric of Eq.~\eqref{eq:chemdist}, the
standard code has total cost $C(\text{id}) = 3{,}581{,}250$.

Of the 190 pairwise swaps, \textbf{42 are strictly improving} (reduce
total cost). The remaining 148 are strictly worsening; no swap is neutral.

Table~\ref{tab:top10} lists the 10 largest-improving swaps. The single
best swap, Arg\,$\leftrightarrow$\,Gln, reduces cost by 367,312
($-10.26\%$). The top eight improving swaps all involve arginine, indicating
that the chemistry-only objective systematically over-penalizes the mutational
neighborhood of the largest 6-fold codon family.

\begin{table}[t]
\centering
\caption{Top 10 improving pairwise swaps under simple chemical distance.
$\Delta C$ is the cost reduction relative to the standard code.}
\label{tab:top10}
\small
\begin{tabular}{@{}llrr@{}}
\toprule
Swap & Degeneracy & $\Delta C$ & \% \\
\midrule
Arg $\leftrightarrow$ Gln & 6$\times$2 & $-$367,312 & $-$10.26 \\
Arg $\leftrightarrow$ Asn & 6$\times$2 & $-$362,696 & $-$10.13 \\
Arg $\leftrightarrow$ Cys & 6$\times$2 & $-$329,794 & $-$9.21 \\
Arg $\leftrightarrow$ His & 6$\times$2 & $-$313,700 & $-$8.76 \\
Arg $\leftrightarrow$ Asp & 6$\times$2 & $-$253,104 & $-$7.07 \\
Arg $\leftrightarrow$ Glu & 6$\times$2 & $-$242,232 & $-$6.76 \\
Arg $\leftrightarrow$ Lys & 6$\times$2 & $-$241,430 & $-$6.74 \\
Arg $\leftrightarrow$ Met & 6$\times$1 & $-$230,034 & $-$6.42 \\
Glu $\leftrightarrow$ Gly & 2$\times$4 & $-$155,500 & $-$4.34 \\
Gly $\leftrightarrow$ Tyr & 4$\times$2 & $-$153,156 & $-$4.28 \\
\bottomrule
\end{tabular}
\end{table}

\subsection{The Improving Swaps Are Not Random}

The 42~improving swaps show strong structural patterns:

\begin{enumerate}
\item \textbf{Arg dominates.} 11 of the 42~improving swaps involve
  arginine, and the top eight improving swaps all involve arginine. The
  chemistry-only objective assigns arginine to an unusually expensive
  mutational neighborhood.

\item \textbf{Cross-degeneracy swaps dominate.} 30 of the 42~improving
  swaps exchange amino acids with different degeneracy counts (for example
  a 6-fold family with a 2-fold or 1-fold family). The metric therefore
  fails to price the structural consequences of degeneracy asymmetry.

\item \textbf{The improving basin is broad.} The top 20 improving swaps
  include exchanges involving 6-fold, 4-fold, 2-fold, and 1-fold families.
  This is not a single isolated anomaly but a persistent feature of the
  chemistry-only landscape.
\end{enumerate}

\subsection{Random Baseline}

To calibrate the standard code's position in the chemical-distance
landscape, we sampled 10,000 degeneracy-preserving random assignments
(shuffling amino acid labels within each degeneracy class independently,
seed = 42). The standard code sits at approximately the \textbf{0.53rd
percentile by cost}: only 53 of the 10,000 sampled random codes are as good
or better. Equivalently, the standard code outperforms 99.47\% of sampled
degeneracy-preserving shuffles under this metric.

This is consistent with the Freeland--Hurst finding~\citep{FreelandHurst1998}
that the standard code outperforms most random codes, but refines it: the
code is globally strong among random alternatives yet still not locally
optimal. The point is therefore not that the code is mediocre under simple
chemistry, but that random-code comparisons alone do not establish that the
right objective has been found.

\subsection{Simulated Annealing}

Simulated annealing with 500,000 iterations (temperature schedule
$T = 10{,}000 \cdot (1/10{,}000)^{t/500{,}000}$, seed = 42) finds codes
with cost 2,709,790, which is 24.33\% cheaper than the standard code.
This is far below the best of the 10,000 sampled random shuffles and
confirms that the chemistry-only objective contains a substantial improving
basin, not merely a few isolated swap anomalies.

\section{Formal Verification}

All results are formalized in Lean~4 (version 4.27.0) using the Mathlib
library. The key modules and theorems are:

\begin{itemize}
\item \texttt{chemicalDistance} in \texttt{AminoAcidZClassification.lean}:
  the distance function of Eq.~\eqref{eq:chemdist}, proved symmetric and
  zero on the diagonal.
\item \texttt{codonMutationWeight} in \texttt{MutationModel.lean}:
  transition/transversion weighting, proved symmetric and zero on equal
  nucleotides.
\item \texttt{totalAssignmentCost} in
  \texttt{CodonAssignmentObjective.lean}: the total cost function of
  Eq.~\eqref{eq:total}.
\item \texttt{chemOnly\_argGln\_improves\_nat} in
  \texttt{ComputableAssignmentCost.lean}:
  a \texttt{native\_decide} proof that the Arg\,$\leftrightarrow$\,Gln
  swap strictly reduces cost under natural-number arithmetic.
\item \texttt{chemOnly\_not\_locally\_optimal} in
  \texttt{CertifiedResults.lean}: the
  theorem that there exist $a \neq b$ with
  $C(\sigma_{a,b}) < C(\text{id})$.
\end{itemize}

The full formalization comprises 49 Lean files totaling over 8,000 lines
of code, with zero axioms and zero unproved assertions (\texttt{sorry}).
Every definition is computable; every inequality is verified by the Lean
kernel via \texttt{native\_decide} on natural-number arithmetic.

\section{Discussion}

\subsection{What the Negative Result Means}

The standard genetic code is not locally optimal under a simple
chemical-distance metric. This does not mean the code is poorly designed.
It means the metric is incomplete. The code was presumably shaped by
constraints that a four-property Euclidean distance cannot capture:

\begin{enumerate}
\item \textbf{Position-dependent mutation effects.} Mutations at the
  second codon position have far larger chemical consequences than
  mutations at the third (wobble) position~\citep{Woese1966}. A metric
  that weights all positions equally conflates these very different
  biological costs.

\item \textbf{Biosynthetic pathway constraints.} Amino acids within the
  same biosynthetic family (e.g., the pyruvate-derived set: Ala, Val, Leu)
  tend to be encoded by neighboring codon boxes. A simple chemical-distance
  metric assigns no special value to this clustering.

\item \textbf{Translation-machinery compatibility.} The aminoacyl-tRNA
  synthetase system imposes structural constraints on which amino acids can
  be assigned to which codon families. A metric that treats all
  reassignments as equally plausible ignores these constraints.

\item \textbf{Charge conservation at wobble.} The wobble position (third
  nucleotide) tolerates synonymous mutations that preserve amino-acid
  identity but can change codon-level properties such as Z-parity. A
  metric that does not distinguish wobble from non-wobble mutations cannot
  capture this asymmetry.
\end{enumerate}

\subsection{Relation to Prior Work}

The Freeland--Hurst result~\citep{FreelandHurst1998} established that the
standard code is better than $\sim$99\% of random codes under a weighted
polar-requirement objective. Our result is complementary, not contradictory.
Under our chemistry-only objective the standard code also beats 99.47\% of
sampled degeneracy-preserving random shuffles, yet still fails local
optimality. The discrepancy arises because rank among random codes is a
global statistic, whereas local optimality asks whether any nearby
perturbation improves the objective. A code can be globally strong relative
to random alternatives and still not be a local minimum of the wrong cost
function.

More recent work by Wnętrzak et al.~\citep{Wnetrzak2018} has shown that
position-dependent weighting significantly changes the optimality landscape.
Our results are consistent with theirs and provide a formal, machine-checked
foundation for the same conclusion.

\subsection{Constructive Implications}

The negative result is constructive: it identifies exactly which swaps
improve the code and by how much. The pattern of improving swaps---dominated
by arginine and by cross-degeneracy exchanges---provides direct clues about
what a correct optimality metric must penalize.

In a companion paper~\citep{Paper2}, we show that extending the metric with
position-dependent weighting and biosynthetic pathway constraints produces
an objective under which the standard code \emph{is} the unique global
optimum over all 363,004 degeneracy-preserving assignments.

\section{Methods}

\subsection{Chemical Properties}

Amino acid properties are drawn from standard reference values: molecular
weights, Kyte--Doolittle hydrophobicity
scale~\citep{KyteDoolittle1982} multiplied by 10 and rounded to integers,
charge at pH~7 ($\{-1, 0, +1\}$), and side-chain volumes from
Zamyatnin~\citep{Zamyatnin1972}. All 20 canonical amino acids are included.

\subsection{Computation}

The optimizer is implemented in Python~3.12 and available in
\texttt{tools/codon\_assignment\_optimizer.py}. All results are
deterministic (random seed = 42 where applicable). Frozen outputs are
stored in \texttt{artifacts/paper1\_frozen\_results.json}. The random
baseline reported in this paper uses 10,000 sampled degeneracy-preserving
shuffles; the simulated annealing run uses 500,000 iterations.

\subsection{Formal Verification}

The Lean~4 formalization is in the \texttt{IndisputableMonolith.Genetics}
namespace of the \texttt{reality} repository. Build command:
\texttt{lake build IndisputableMonolith.Genetics}. The build requires
Lean~4.27.0 and a pre-downloaded Mathlib cache (\texttt{lake exe cache get}).

\section{Conclusion}

The standard genetic code is not locally optimal under a simple
chemical-distance metric. Forty-two of 190 pairwise amino-acid swaps
strictly reduce the total weighted mutation cost, with the best swap
(Arg\,$\leftrightarrow$\,Gln) improving the objective by 10.26\%.
At the same time, the standard code still outperforms 99.47\% of sampled
degeneracy-preserving random shuffles. This combination of facts is
machine-checked in Lean~4, eliminating the possibility of computational
error or off-by-one mistakes in the enumeration.

The finding does not diminish the code's remarkable structure. Instead, it
sharpens the question: what objective \emph{does} the code minimize? The
answer must include position-dependent mutation effects, biosynthetic
pathway constraints, and translation-machinery compatibility---all of which
are absent from the simple chemical-distance metric tested here.

\section*{Data Availability}

All code, data, and Lean source files are available in the
\texttt{reality} repository. The optimizer script, frozen result tables,
and Lean formalization can be used to reproduce every number in this paper.

\section*{Acknowledgments}

The Lean formalization uses the Mathlib library maintained by the
\texttt{leanprover-community}. Computational resources were standard
desktop hardware.

\begin{thebibliography}{99}
\bibitem[Crick(1968)]{Crick1968}
Crick, F.H.C. (1968). The origin of the genetic code.
\textit{J.~Mol.~Biol.} \textbf{38}(3), 367--379.

\bibitem[Freeland and Hurst(1998)]{FreelandHurst1998}
Freeland, S.J. and Hurst, L.D. (1998). The genetic code is one in a
million. \textit{J.~Mol.~Evol.} \textbf{47}(3), 238--248.

\bibitem[Grantham(1974)]{Grantham1974}
Grantham, R. (1974). Amino acid difference formula to help explain protein
evolution. \textit{Science} \textbf{185}(4154), 862--864.

\bibitem[Haig and Hurst(1991)]{HaigHurst1991}
Haig, D. and Hurst, L.D. (1991). A quantitative measure of error
minimization in the genetic code. \textit{J.~Mol.~Evol.} \textbf{33}(5),
412--417.

\bibitem[Kyte and Doolittle(1982)]{KyteDoolittle1982}
Kyte, J. and Doolittle, R.F. (1982). A simple method for displaying the
hydropathic character of a protein. \textit{J.~Mol.~Biol.} \textbf{157}(1),
105--132.

\bibitem[Wakeley(1996)]{Wakeley1996}
Wakeley, J. (1996). The excess of transitions among nucleotide
substitutions: new methods of estimating transition bias underscore its
significance. \textit{Trends Ecol.~Evol.} \textbf{11}(4), 158--163.

\bibitem[Wnętrzak et~al.(2018)]{Wnetrzak2018}
Wnętrzak, M., Błażej, P., Mackiewicz, D. and Mackiewicz, P. (2018).
The optimality of the standard genetic code assessed by an eight-quality
model. \textit{BMC Evol.~Biol.} \textbf{18}(1), 71.

\bibitem[Woese(1965)]{Woese1965}
Woese, C.R. (1965). On the evolution of the genetic code.
\textit{Proc.~Natl.~Acad.~Sci.~USA} \textbf{54}(6), 1546--1552.

\bibitem[Woese et~al.(1966)]{Woese1966}
Woese, C.R., Dugre, D.H., Dugre, S.A., Kondo, M. and Saxinger, W.C.
(1966). On the fundamental nature and evolution of the genetic code.
\textit{Cold Spring Harb.~Symp.~Quant.~Biol.} \textbf{31}, 723--736.

\bibitem[Zamyatnin(1972)]{Zamyatnin1972}
Zamyatnin, A.A. (1972). Protein volume in solution. \textit{Prog.~Biophys.
Mol.~Biol.} \textbf{24}, 107--123.

\bibitem[Paper~2]{Paper2}
Washburn, J. (2026). The standard genetic code is globally optimal under
position-weighted mutation cost. \textit{In preparation}.
\end{thebibliography}

\end{document}
